\begin{abstract}
This report compares clustered scanning, Bitplanes and Backward Walk on one collection of
binary documents. The calculations explain how cluster size, bit splits and key lookups
affect work and memory. Experiments check those estimates and measure when each method
is faster under the same recall and RAM requirement.
\end{abstract}

\section{Introduction}
\subsection{Exa method and end to end pipeline}
Exa describes a search pipeline that shortens Matryoshka embeddings, stores document values
as binary signs, selects nearby clusters, and scores their documents with the floating-point
query~\cite{exa}. The results can then be reranked with uncompressed vectors. We refer to that
approach as the original method, or A. The experiments use our implementation of this search
stage; they do not measure Exa's production system.

Documents are indexed before queries arrive. At query time, A opens some clusters and scores
their documents. The blue box is the step that B and C change.

\begin{figure}[H]\centering
\begin{tikzpicture}[node distance=6mm and 6mm]
\node[box,text width=25mm] (text) {Document text};
\node[box,text width=32mm,right=of text] (embed) {Embed and retain\\the first $d$ values};
\node[box,text width=27mm,right=of embed] (bits) {Pack document\\signs};
\node[box,text width=25mm,right=of bits] (store) {Assign clusters\\and build index};
\draw[arr] (text)--(embed);\draw[arr] (embed)--(bits);\draw[arr] (bits)--(store);
\node[box,text width=22mm,below=8mm of text] (q) {Query text};
\node[box,text width=25mm,right=3mm of q] (qvec) {Encode\\query};
\node[box,text width=28mm,right=3mm of qvec] (route) {Select $P$ of $C$\\clusters};
\node[scan,text width=25mm,right=3mm of route] (local) {Filter and\\score (A)};
\node[box,text width=22mm,right=3mm of local] (out) {Rerank\\results};
\draw[arr] (q)--(qvec);\draw[arr] (qvec)--(route);\draw[arr] (route)--(local);\draw[arr] (local)--(out);
\end{tikzpicture}
\caption{Exa's described end-to-end flow. B and C replace the blue search step. The benchmark
starts from cached embeddings and times cluster selection and search.}
\end{figure}

\subsection{Idea background}
The query itself tells us the best possible combination of document signs. To follow the
arithmetic, consider a five-value example:
\[
q=[0.7,\ 0.5,\ -0.4,\ -0.6,\ 0.3].
\]
A document bit of 1 represents $+1$; a bit of 0 represents $-1$. At the first position,
$0.7\times(+1)=0.7$ is better than $0.7\times(-1)=-0.7$. At the third position,
$-0.4\times(-1)=0.4$ is better than $-0.4\times(+1)=-0.4$. Following this rule at every
position gives \texttt{11001}, with score $0.7+0.5+0.4+0.6+0.3=2.5$.

A mismatch reduces the score by twice the size of that query value. Mismatching the first
position loses 1.4; mismatching the last loses 0.6. In general,
\begin{equation}
 S(q,b)=Q-2p,\qquad Q=\sum_i|q_i|,\qquad
 p=\sum_{i:\,b_i\ne\operatorname{sign}(q_i)}|q_i|.
 \label{eq:score}
\end{equation}
Here $p$ adds the sizes of the query values at mismatched positions. For example,
\texttt{11000} differs only at the last bit, giving score $2.5-2(0.3)=1.9$.

The question is whether an index can find documents near the preferred pattern without
scoring every document. Bitplanes use the query values to choose which bits to split on.
Backward walk starts with a short preferred pattern and tries other patterns when needed.
The preferred pattern may not exist in the collection, so both methods must retain alternatives.

\newpage
\subsection{Exploring two new methods: ``Bitplanes'' and ``Backward Walk''}
\subsubsection{Bitplanes}
Method B keeps the packed rows and adds a second arrangement of their signs. Each
\emph{bitplane} stores one bit position across all documents in a cluster. The figure shows
three documents with the query from Section 1.2. A candidate mask records which local rows remain.

\begin{figure}[H]\centering
\begin{tikzpicture}
\node[box,text width=35mm] (route) at (0,0) {Select clusters\\as in A};
\node[branch,text width=43mm] (branch) at (4.8,0) {Read bitplanes\\and split candidates};
\node[box,text width=35mm] (score) at (9.6,0) {Score surviving rows\\same full score};
\draw[arr] (route)--(branch);\draw[arr] (branch)--(score);
\node[align=left,font=\small] at (1.25,-2.1) {Packed rows\\ID 2: \texttt{10100}\\ID 3: \texttt{11001}\\ID 4: \texttt{11000}};
\draw[arr] (3,-2.1)--node[above,font=\small]{store columns}(4.25,-2.1);
\node[align=center,font=\small] at (7,-2.1) {\begin{tabular}{c|ccc}
 & ID 2 & ID 3 & ID 4\\\hline
Bit 1 &1&1&1\\
\color{Branch}Bit 2 &\color{Branch}0&\color{Branch}1&\color{Branch}1\\
Bit 3 &1&0&0\\
Bit 4 &0&0&0\\
Bit 5 &0&1&0
\end{tabular}};
\node[branch,text width=38mm] (parent) at (5,-4.1) {Parent mask \texttt{111}};
\node[branch,text width=37mm] (yes) at (2.3,-5.3) {Match bit 2: \texttt{011}\\IDs 3,4; penalty 0};
\node[box,text width=37mm] (no) at (7.7,-5.3) {Other: \texttt{100}\\ID 2; penalty 0.5};
\draw[arr] (parent)--(yes);\draw[arr] (parent)--(no);
\end{tikzpicture}
\caption{Orange marks the changed stage and the selected plane. Both child masks retain all
three positions. On a larger cluster, this physical width is a major cost.}
\label{fig:bitplanes}
\end{figure}

The absolute sizes of the query values give this position order: 1, 4, 2, 3, 5. For a positive query
value, a parent mask $A$ is split into $A\mathbin{\&}B_i$ and
$A\mathbin{\&}\neg B_i$. The first child matches; the other adds $|q_i|$ to its penalty.
Negative query values reverse the roles. Alternatives remain in a priority queue,
ordered by known penalty, because a preferred early bit does not determine the final ranking.

Equation~\ref{eq:score} gives a safe upper bound: a branch with known penalty $p$ can score at
most $Q-2p$. If IDs 3 and 4 score 2.5 and 1.9, the alternative with penalty 0.5 has bound 1.5 and
cannot improve the top two. Equal-score cases retain the common document-ID tie rule.

\begin{lstlisting}
bitplanes(query, clusters, K, leaf_size, node_budget):
    table, best = prepare_score_table(query), empty_top_k(K)
    order = positions_by_decreasing_abs_value(query)
    queue = full candidate mask for each selected cluster
    while queue is not empty:
        node = remove_lowest_penalty(queue)
        if cannot_improve(node, query, best): continue
        if node_budget_is_reached(): break
        if node.count <= leaf_size or no_bits_remain(node):
            score_set_documents(table, node.mask, best)
        else:
            queue_nonempty_children(split(node, order))
    return best.sorted_by_score_then_id()
\end{lstlisting}
The leaf size decides when to stop splitting and score the remaining rows. Unlimited search
with safe bounds is exact within opened clusters; a finite work limit can lose results.

\newpage
\subsubsection{Backward Walk}
Method C forms a short key from a chosen set of $h$ bit positions. In each cluster, documents
with the same key are stored consecutively. A directory maps an occupied key to the start
and length of that row range. Search begins at the query's preferred key and enumerates
alternatives in increasing weighted mismatch penalty.

\begin{figure}[H]\centering
\begin{tikzpicture}
\node[box,text width=32mm] (route) at (0,0) {Select clusters\\as in A};
\node[keys,text width=45mm] (walk) at (4.7,0) {Enumerate weighted keys\\and read the directory};
\node[box,text width=32mm] (score) at (9.4,0) {Score listed rows\\same full score};
\draw[arr](route)--(walk);\draw[arr](walk)--(score);
\node[keys,text width=25mm] (query) at (-.3,-2.2) {Query key\\\texttt{10}\\bits 2,3};
\node[align=center,font=\small] (directory) at (4.8,-2.2) {
\begin{tabular}{cccl}\toprule
Key&Penalty&Cluster 0&Cluster 1\\\midrule
\texttt{10}&0&[0]&[3,4]\\
\texttt{11}&0.4&[]&[]\\
\texttt{00}&0.5&[5]&[]\\
\texttt{01}&0.9&[1]&[2]\\\bottomrule
\end{tabular}};
\draw[arr](query)--(directory);
\node[keys,text width=28mm] at (10.5,-2.2) {Stored entry:\\key, start, count\\$\downarrow$\\contiguous rows};
\node[keys,text width=24mm] (first) at (0,-4.1) {\texttt{10}\\score 3 rows};
\node[box,text width=25mm] (second) at (3.5,-4.1) {\texttt{11}\\bound 1.7};
\node[box,text width=25mm] (third) at (7,-4.1) {\texttt{00}\\bound 1.5};
\node[box,text width=25mm] (fourth) at (10.5,-4.1) {\texttt{01}\\bound 0.7};
\draw[arr](first)--(second);\draw[arr](second)--(third);\draw[arr](third)--(fourth);
\end{tikzpicture}
\caption{Purple marks the changed stage. Here the first key returns IDs 0,3,4. Their top two
scores are 2.5 and 1.9, so the next key's bound 1.7 already permits stopping. The remaining
entries show the order that a query requiring more work would follow.}
\label{fig:keys}
\end{figure}

The ordering uses the sum of absolute query values at the flipped positions. Counting flipped bits alone
would give the wrong order when query weights differ. There is one shared enumeration queue
across opened clusters, followed by a lookup in each cluster's directory. Empty lookups also
take time. Removing a key constraint is not a single ordinary hash lookup: the newly admitted
patterns must be enumerated or supported by a different data structure.

\begin{lstlisting}
backward_walk(query, clusters, K, key_positions, limits):
    table, best = prepare_score_table(query), empty_top_k(K)
    keys = unique_keys_by_increasing_penalty(query, key_positions)
    for key, penalty in keys:
        if cannot_improve(penalty, query, best): break
        for cluster in clusters:
            if work_limit_reached(limits):
                return best.sorted_by_score_then_id()
            rows = cluster.directory.lookup(key)
            score_each_row(table, rows, best)
    return best.sorted_by_score_then_id()
\end{lstlisting}

A short key omits some bit positions, so every retrieved row still receives the complete
score. Finding enough candidates does not by itself prove sufficient recall. Exhausting the
keys or ruling out every remaining bound is exact within opened clusters; stopping at a
candidate or lookup budget is approximate. Neither B nor C can recover a document excluded
by the initial routing.
